\section{Current Law of the Latin Church}

\subsection{The governing obligation}

CIC canon 277 \S1 obliges clerics to perfect and perpetual continence for the sake of the kingdom and therefore to celibacy, which the canon commends as a special gift permitting closer adherence to Christ and freer service of God and humanity. Sections 2 and 3 require appropriate prudence in relationships and authorize the diocesan bishop to issue more specific norms and judge particular cases. The canon is juridic, ascetical, and relational; it does not reduce fidelity to avoiding one prohibited act.

A man becomes a cleric through diaconal ordination (can. 266 \S1). A candidate for the presbyterate and an unmarried candidate for the permanent diaconate must publicly assume celibacy before diaconal ordination unless already bound by perpetual religious profession (can. 1037). The corresponding commitment is not ordinarily a diocesan priest's religious vow. Canon 276 places it within the wider clerical vocation to holiness, prayer, Eucharistic life, Scripture, retreat, and spiritual means.

\subsection{Selection for presbyterate and episcopate}

The ordinary Latin discipline selects presbyteral candidates who freely undertake celibacy. The law admits limited cases of married men ordained by authority of the Apostolic See, especially some former non-Catholic ministers and cases governed by the personal-ordinariate norms. These are authorized applications of Latin law, not Easternization and not proof that the norm has disappeared. The governing rescript or special norm matters in each case.

Latin bishops are ordinarily chosen from celibate presbyters. This selection does not imply that marriage is incompatible with priestly character; it continues the ancient practice of choosing bishops not bound by marriage.

\subsection{Orders prevents a later marriage}

CIC canon 1087 makes sacred orders a diriment impediment: one in orders invalidly attempts marriage. Canon 1394 \S1 separately attaches \emph{latae sententiae} suspension to a cleric attempting marriage, even only civilly, with progressively graver penalties if he does not reform after warning or continues scandal, up to dismissal from the clerical state. Invalidity, moral culpability, penal imputability, and procedural consequence are separate findings.

Canon 1042, 1\textsuperscript{o}, simply impedes a married man from receiving orders unless he is legitimately destined to the permanent diaconate. Calling the ordination of a married permanent deacon a ``dispensation from celibacy'' is inaccurate: the code itself establishes his route. A different married candidate requires the competent dispensation or special law.

\subsection{Loss of state does not undo ordination or automatically permit marriage}

Sacred ordination, once validly received, never becomes invalid (can. 290). Loss of the clerical state removes the rights and obligations attached to clerical status as the law specifies, but canon 291 expressly provides that it does not itself dispense the obligation of celibacy; that dispensation is reserved to the Roman Pontiff. ``Laicization'' is therefore not sacramental erasure and is not, standing alone, permission to marry.

\subsection{Chastity, delict, and penalty}

Not every sin against chastity is automatically a canonical delict, and not every delict carries an automatic penalty. Current canon 1395 addresses concubinage, persistent external sexual sin with scandal, other public sexual delicts, and sexual wrongdoing committed by force, threats, abuse of authority, or coercion. Canon 1398 specifies offenses involving minors and persons whom the law equally protects, including enumerated pornographic conduct. Canon 1394 addresses attempted marriage. Penal typicity, external violation, imputability, proof, prescription, procedure, and the rights of all parties must be applied.

The revised Book VI, promulgated by \emph{Pascite gregem Dei} and effective 8 December 2021, controls this penal description. Civil reporting duties, protective measures, and the welfare of victims are not displaced by internal canonical categories. No reader should infer from an act's failure to fit one penal canon that the act was morally permissible or pastorally safe.

\begin{threestable}{Question}{General Latin answer}{Fact that can change the analysis}
May a Latin priest marry after ordination? & No; sacred orders invalidates the attempt, and an attempted marriage also has penal consequences. & A Roman dispensation joined to loss of clerical state can address the impediment and celibacy obligation.\\
May a married man become a priest? & Not under the ordinary formation route; limited Holy See-authorized paths exist. & Prior non-Catholic ministry, ordinariate law, individual rescript, validity of marriage, and spouse/family circumstances.\\
May a married man become a deacon? & Yes, if legitimately destined to the permanent diaconate and all requirements are met. & Minimum age, wife's written consent, suitability, family stability, conference norms, and impediments.\\
May an unmarried permanent deacon marry? & No; he publicly assumed celibacy and is in sacred orders. & Only a competent Apostolic See dispensation can remove the impediment where available.\\
Does dismissal permit marriage? & Not by itself. & The papal rescript must be read for both loss of state and dispensation.\\
Does every sexual sin incur a canonical penalty? & No. & The exact external conduct, applicable penal canon or precept, imputability, evidence, and process.\\
\end{threestable}
